% =============================================================================
%                    TEMPLATE VOOR CURSUSSAMENVATTING
%                    KU Leuven - School Macros v4.0
% =============================================================================
% Dit is een starttemplate voor nieuwe documenten.
% Kopieer dit bestand en pas het aan voor je eigen cursus.
%
% Tip: Compileer TWEEMAAL voor correcte formularium en symbolenlijst!
% =============================================================================

\documentclass[a4paper,11pt]{article}
\usepackage{school-macros}
\enablecompactmode

% v4.0 Optioneel: Compact Mode voor cheatsheets
% \enablecompactmode

% =============================================================================
%                           DOCUMENT INFO
% =============================================================================
% Pas deze gegevens aan voor je cursus:
\newcommand{\vaknaam}{Object gericht-programmeren}
\newcommand{\auteur}{Ruben Ryckaert}

\title{\vaknaam{} --- Synthax Java}
\author{\auteur}

\date{\today}


% =============================================================================
%                              DOCUMENT
% =============================================================================
\begin{document}
\maketitle

 
\begin{abstract}
Deze samenvattingen zijn beschikbaar op GitHub waar je kunt bijdragen aan verbeteringen: \url{https://github.com/Eggmansmile/Samenvattingen-Ku-leuven-Industriele-ingenieurs}

Je bent helemaal vrij en ik raad je zelf aan om info toe te voegen aan de documenten als je de uitleg niet goed vindt.
Je leert je vak en je raakt dan ook bekend met LaTeX wat handig is voor later.
\end{abstract}

\vspace{1cm}


Je mag voor het examen een \textbf{handgeschreven 
4-pagina's A4-formularium meenemen} met daarop alles wat je nodig hebt.
Dit document moet je tonen wat de belangrijkste dingen zijn die je moet weten voor het examen.
Ik raad ja aan om dit document te lezen en de belangrijkste dingen over te schrijven op je formularium.
Je hebt ook op studforum grote hoeveelheden oefenexamens staan die je zeker moet maken. 
\url{https://cloud.studforum.net/apps/files/?dir=/GTS-Fase2/Semester%201/Courses%20After%202021/Nederlandse%20Groep/Objectgerichte%20Softwareontwikkeling&fileid=794166}
\newline 

W3schools heeft een goede Java tutorial: \url{https://www.w3schools.com/java/}
Bekijk die zeker als je niet weet wat alles doet. Dit is een goed beginpunt.


\section{Basis java}
\subsection{Import statements}
Om klassen van andere packages te gebruiken, 
moeten deze eerst geïmporteerd worden met een import statement.\\

\begin{codeblock}[Import utils]{java}
    import java.util.*; // importeert alle klassen van de package java.util
    import java.util.ArrayList; // importeert enkel de ArrayList klasse
\end{codeblock}

\subsection{Typische Java klasse structuur \& waarden}
Een typische Java klasse heeft de volgende structuur:
\begin{codeblock}[Hoe java doc starten?]{java}
    import java.io.*; //zoals File, FileReader, BufferedReader, etc

    import java.util.*; //zoals ArrayList, Scanner, Random, etc
    

    public class KlasseNaam  // klasse definitie
    {
        // opstellen van variabelen (attributen)
        private int attribuut1;
        private String attribuut2;
        private float attribuut3; // niet alle attributen moeten in de constructor zitten
        public int attribuut4; // kan ook public zijn 

        // constructor (Je creeert een instance van een klasse met een constructor)
        public KlasseNaam(int attribuut1, String attribuut2) 
        {
            this.attribuut1 = attribuut1;
            this.attribuut2 = attribuut2;
        }

        // methoden
        public void methode1() 
        {
            // methode implementatie
        }

        public int methode2(String parameter) 
        {
            // methode implementatie
            return 0;
        }

        #alle soorten characters:
        int #gehele getallen
        int [] #single dim array
        int [] [] #multi dim arrays

        String #tekst
        String [] #array van teksten

        float #decimale getallen

        boolean #true of false

        char #enkel character

        double #decimale getallen (anders dan float)
        double [] #array van double

        //belangrijke na elke lijn ;
    }
\end{codeblock}

\textbf{Public} betekent dat andere klassen deze variabele of methode
 kunnen gebruiken.
\textbf{Private} betekent dat alleen deze klasse deze variabele of methode 
kan gebruiken.

\subsection{Methoden}

een methode zit in je klasse. Hiermee
kun je functionaliteit toevoegen aan je klasse.

\begin{codeblock}[methode voorbeeld]{java}

    public int #//*return type*// methodeOptelling(int parameter1, int parameter2) 
    {  
        int resultaat = parameter1 + parameter2;
        return resultaat; // geeft het resultaat terug

    }

    //een methode met return type 'void' geeft niets terug:
    public void methodePrint(String tekst) 
    {  
        System.out.println(tekst); // print de tekst naar de console
    }

    Je kunt alle characters returnen die je wilt.
\end{codeblock}

\section{If, for en while}
\subsection{If statements}
\begin{codeblock}[if statement]{java}
    if (voorwaarde) {
        // code die uitgevoerd wordt als de voorwaarde waar is
    } else if (andereVoorwaarde) 
    {
        // code die uitgevoerd wordt als de andere voorwaarde waar is
    } else 
    {
        // code die uitgevoerd wordt als geen van de voorwaarden waar is
    }

    if en else zijn een sort scan, dus eerst wordt if gedaan en als
    dat niet waar is, gaat die naar else if en zo verder.
    Moest je nu gewoon if statements gebruiken, dan worden ze allemaal
    uitgevoerd.

    Dit zijn belangrijke voorwaarden:
    ==  // gelijk aan
    !=  // niet gelijk aan
    >   // groter dan
    <   // kleiner dan
    >=  // groter dan of gelijk aan
    <=  // kleiner dan of gelijk aan
    &&  // en
    ||  // of
\end{codeblock}

\subsection{For loops}
\begin{codeblock}[for loop]{java}
    for (init; voorwaarde; increment) 
    {
        // code die herhaald wordt
    }

    // voorbeeld: print getallen van 0 tot 9
    for (int i = 0; i < 10; i++) // i begint bij 0, zolang i kleiner is dan 10, verhoog i met 1
    //vergeet de ; niet tussen de statements (init, voorwaarde, increment)
    {
        System.out.println(i); //output is 0 1 2 3 4 5 6 7 8 9
    }
\end{codeblock}

\subsection{While loops}
While loops zijn gelijkaardig aan for loops, maar ze zijn
meer geschikt voor situaties waarin je niet weet hoeveel iteraties je nodig hebt.
\begin{codeblock}[while loop]{java}
    int i = 0; // initialisatie
    while (i < voorwaarde) // zolang i kleiner is dan de voorwaarde
     {
        // code die herhaald wordt zolang de voorwaarde waar is
        i++; // vergeet niet i aan te passen om een oneindige loop te vermijden
        i--; // of i verlagen
    }
\end{codeblock}


\subsection{Arrays}

een array is wiskundig een vector.
een array van strings zijn gewoon meerdere teksten in 1 variabele.
\begin{codeblock}[array]{java}
    // array van gehele getallen met 5 elementen
    int[] getallen = new int[5];

    // array van strings met 3 elementen
    String[] namen = new String[3]; // een array zijn lengte is vast. pas op want dit is niet de index.

    // array initialisatie met waarden
    int[] cijfers = {90, 85, 78, 92, 88};

    // toegang tot array elementen
    int eersteGetal = getallen[0]; // eerste element (index 0) 
    String eersteNaam = namen[0]; // eerste element (index 0)

    //belangrijke array functies
    int lengte = cijfers.length; // lengte van de array
    int laatsteIndex = cijfers.length - 1; // laatste index van de array
    int middelsteIndex = cijfers.length / 2; // middelste index van de array
    int 
    // itereren over een array
    for (int i = 0; i < cijfers.length; i++) 
    {
        System.out.println(cijfers[i]);
    }

    je hebt nog een foreach loop. Deze moet je zeker kennen voor ArrayLists maar
    is ook hier handig.

    for (int cijfer : cijfers) //zorg dat je de : gebruikt in plaats van ;
    // je zegt dus voor elk cijfer (die gebruik je in de loop) in cijfers (de array)
    {
        System.out.println(cijfer);
    }
\end{codeblock}


\subsection{Multi dim arrays}

\begin{codeblock}[multi dim array]{java}
    // 2D array (matrix) van gehele getallen met 3 rijen en 4 kolommen
    int[][] matrix = new int[3][4];

    // initialisatie van een 2D array
    int[][] tabel = {
        {1, 2, 3, 4},
        {5, 6, 7, 8},
        {9, 10, 11, 12}
    };

    // toegang tot elementen in een 2D array
    int element = tabel[1][2]; // element in de tweede rij en derde kolom (waarde is 7)

    // itereren over een 2D array
    for (int i: tabel) //voor elke rij in tabel
    {
        for (int j: matrix[i]) //voor elke kolom in die rij
        {
            System.out.println(matrix[i][j]);
            //output: 1 2 3 4 5 6 7 8 9 10 11 12
        }

        // je gaat eerst over elke element in de rij en 
        dan naar de volgende rij.
    }
\end{codeblock}

\subsection{ArrayLists}
een opstelling van een ArrayList:

\begin{codeblock}[arraylist]{java}
    import java.util.ArrayList; //importeer de ArrayList klasse

    // maak een ArrayList van Strings

    private ArrayList<String> namen;

    in constructor:
        public MijnKlasse() 
        {
            namen = new ArrayList<String>();
        }

    // voeg elementen toe aan de ArrayList
    namen.add("Alice");
    namen.add("Bob");
    namen.add("Charlie");

    // toegang tot elementen in de ArrayList
    String eersteNaam = namen.get(0); // eerste element (index 0)

    // verwijder een element uit de ArrayList
    namen.remove(1); // verwijdert het element op index 1 ("Bob")
    namen.remove("Bob"); // verwijdert het element "Bob" als het bestaat

    // grootte van de ArrayList
    int grootte = namen.size(); // aantal elementen in de ArrayList

    // itereren over een ArrayList
    for (String naam : namen) 
    {
        System.out.println(naam);
        //output: Alice Charlie
    }

    //arraylists kunnen ook instancies bevatten van
    andere klassen,
    bijvoorbeeld een arraylist van personen:

    private ArrayList<Persoon> personen;

    in constructor:
        public MijnKlasse() 
        {
            personen = new ArrayList<Persoon>();
        }

    // voeg een persoon toe aan de ArrayList
    public void voegPersoonToe(String naam, int leeftijd) 
    {
    personen.add(new Persoon(naam, leeftijd)); //persoon(constructor)
    }

    Je kunt dan ook methoden van de Persoon klasse gebruiken:
    //scan die personen boven een bepaalde leeftijd verwijderd 
    uit de lijst en die teruggeeft
    public ArrayList<Persoon> verwijderOuderePersonen(int leeftijdsGrens) 
    {
        ArrayList<Persoon> verwijderdePersonen = new ArrayList<Persoon>();

        for (Persoon persoon : personen) 
        {
            //methode van Persoon klasse
            if (persoon.getLeeftijd() > leeftijdsGrens) 
            {
                verwijderdePersonen.add(persoon);
                personen.remove(persoon);
            }
        }
    }
\end{codeblock}


\section{Class extends}
In Java kun je een nieuwe klasse maken die de eigenschappen en methoden van een bestaande klasse overneemt door middel van het sleutelwoord \texttt{extends}. Dit wordt ook wel \textit{inheritance} genoemd.

\begin{codeblock}[class extends]{java}
    // Bestaande klasse
    public class Dier 
    {

    private String naam;
    private int leeftijd;
    private String locatie;

    public Dier(String naam, int leeftijd, String locatie) {
        this.naam = naam;
        this.leeftijd = leeftijd;
        this.locatie = locatie;
        }

        public void getInfo() {
            System.out.println("Naam: " + naam + ", Leeftijd: " + leeftijd + ", Locatie: " + locatie);
        }
    }

    // Nieuwe klasse die Dier uitbreidt
    public class Hond extends Dier 
    {

        private String ras;

        public Hond(String naam, int leeftijd, String locatie, String ras) {
            super(naam, leeftijd, locatie); // roept de constructor van de superklasse Dier aan
            this.ras = ras;
        }
        // je moet methodes niet opnieuw schrijven want een hond is ook een dier
        public void maakGeluid() {
            System.out.println("Hond blaft");
        }

        //alleen hond kan deze methode gebruiken
        public void speelFetch() {
            System.out.println("Hond speelt fetch");
        }
    }

    // Gebruik van de klassen
    public class Main 
    {
        ArrayList<Dier> dieren = new ArrayList<Dier>();
        dieren.add(new Dier("Leeuw", 5, "Afrika"));
        dieren.add(new Hond("Buddy", 3, "Thuis", "Labrador"));
        //je kunt een hond ook toevoegen omdat een hond een dier is

        for (Dier dier : dieren) {
            dier.getInfo(); // werkt voor zowel Dier als Hond
            //dier.maakGeluid(); // Fout: Dier heeft deze methode niet
            if (dier instanceof Hond) {
                ((Hond) dier).maakGeluid(); // Correct: cast naar Hond om de methode aan te roepen
            }
        }
    }


\section{Hashmaps} 

% TODO: #62 dit kwam op het examen maar geen idee wat het is? Iemand uitleggen?

\section{Scanner}

De code om van de Scanner krijg je op het examen maar niet hoe je die gebruikt.
\begin{codeblock}[scanner]{java}
    import java.util.Scanner; //importeer de Scanner klasse

    public class MijnKlasse 
    {
    
    public MijnKlasse() {
    }

    public int sumOfTwoNumbers(int a, int b) {
        return a + b;
    }

public void readFromFile(String filename) {
  try {
    Scanner sc = new Scanner(new File(filename));
    while(sc.hasNext()) 
    {
        String line = sc.nextLine();
        System.out.println(line);
    }
    of met integers:
    while(sc.hasNextInt())
    {
        int a = sc.nextInt();
        int b = sc.nextInt();
        int resultaat = sumOfTwoNumbers(a, b);
        System.out.println(resultaat);
    }
    sc.close();
  }
  catch (FileNotFoundException fnfe) 
  {
    System.out.println("file not found");
    }
    }
}

//belangrijke methodes van Scanner:
hasNext(): controleert of er nog meer tokens zijn
next(): leest het volgende token als een String
nextInt(): leest het volgende token als een int
nextLine(): leest de volgende volledige regel als een String
nextFloat(): leest het volgende token als een float
nextDouble(): leest het volgende token als een double
\end{codeblock}

\section{Random}

om willekeurige getallen te genereren in Java, kun je de \texttt{Random} klasse uit de \texttt{java.util} package gebruiken.

\begin{codeblock}[random]{java}
    import java.util.Random; //importeer de Random klasse

    public class MijnKlasse 
    {
    
    public MijnKlasse() {
    }

    public void generateRandomNumbers() {
        Random rand = new Random(); // maak een Random object

        // genereer een willekeurig geheel getal tussen 0 (inclusief) en 100 (exclusief)
        int randomInt = rand.nextInt(100);
        System.out.println("Willekeurig geheel getal: " + randomInt);

        // genereer een willekeurig decimaal getal tussen 0.0 (inclusief) en 1.0 (exclusief)
        float randomFloat = rand.nextFloat();
        System.out.println("Willekeurig decimaal getal: " + randomFloat);

        // genereer een willekeurig geheel getal binnen een specifiek bereik, bijvoorbeeld tussen 50 en 150
        int min = 50;
        int max = 150;
        int randomInRange = rand.nextInt(max - min) + min;
        System.out.println("Willekeurig geheel getal tussen " + min + " en " + max + ": " + randomInRange);
    }
}
\end{codeblock}

\section{UML diagrammen}
% TODO: #63 UML diagrammen uitleggen 

\section{Hashmaps} 

% TODO: #62 dit kwam op het examen maar geen idee wat het is? Iemand uitleggen?

\end{document}
